\chapter{Conclusão}

% É a parte que sintetiza os argumentos e elementos contidos no desenvolvimento do trabalho, em que são apresentadas as conclusões próprias da pesquisa, retomando o problema inicial e os objetivos e revendo as principais contribuições do estudo.

O desenvolvimento do OmniBot representou a consolidação de uma plataforma robótica móvel omnidirecional voltada ao ensino e à pesquisa em Engenharia Mecânica, integrando conceitos de robótica móvel, Internet das Coisas (IoT) e controle distribuído via ROS 2. A proposta atendeu plenamente aos objetivos definidos, ao demonstrar a viabilidade de construir um robô modular e funcional, capaz de operar em conjunto com sistemas de software livres e escaláveis.

A adoção do ESP32 DevKit V1 como unidade de processamento e da BitDogLab como interface de controle remoto revelou-se eficaz, comprovando que tecnologias acessíveis podem ser aplicadas de forma robusta em ambientes acadêmicos. A arquitetura baseada em micro-ROS e HTTP garantiu flexibilidade de integração com outros dispositivos e sistemas em rede, tornando o OmniBot uma base para futuras aplicações educacionais e laboratoriais.

No âmbito educacional, o OmniBot se destacou como ferramenta didática de grande potencial, capaz de auxiliar na aplicação prática de conteúdos previstos no Projeto Pedagógico do Curso (PPC) de Engenharia Mecânica do IFPI, especialmente nas disciplinas de Robótica e Automação. A plataforma permite que estudantes explorem, de maneira prática, os princípios de cinemática, controle, sensoriamento e comunicação em rede, promovendo um aprendizado experimental alinhado às demandas da indústria.

Os testes realizados no LABIRAS confirmaram o pleno funcionamento do sistema em todos os modos de operação. O robô apresentou movimentação adequada, graças à matriz cinemática das rodas \textit{mecanum}, e respondeu de forma estável aos comandos remotos via BitDogLab e ROS 2. O chassi impresso em 3D, projetado no Autodesk Fusion 360, demonstrou boa resistência mecânica, modularidade e facilidade de manutenção, validando o uso da manufatura aditiva como processo eficiente para prototipagem de plataformas robóticas educacionais.

Além do sucesso técnico e funcional, o OmniBot reforça a importância de projetos integradores no contexto da formação em Engenharia Mecânica, unindo aspectos de projeto mecânico, eletrônica embarcada, programação e automação industrial em um único sistema. O caráter aberto e documentado do projeto estimula sua replicação e evolução por outros estudantes e pesquisadores, ampliando seu impacto acadêmico e científico.